% Generated draft: requires mathematical and source review.
\documentclass[11pt]{article}
\usepackage[margin=25mm]{geometry}
\usepackage{fontspec}
\setmainfont{Latin Modern Roman}
\usepackage{amsmath,amssymb,mathtools,unicode-math}
\setmathfont{Latin Modern Math}
\usepackage{enumitem,needspace,xurl,hyperref}
\setlength{\parindent}{0pt}
\setlength{\parskip}{5pt}
\setlength{\emergencystretch}{3em}
\newcommand{\N}{\mathbb{N}}
\newcommand{\Z}{\mathbb{Z}}
\newcommand{\Q}{\mathbb{Q}}
\newcommand{\R}{\mathbb{R}}
\newcommand{\C}{\mathbb{C}}
\newcommand{\F}{\mathbb{F}}
\newcommand{\A}{\mathbb{A}}
\newcommand{\PP}{\mathbb{P}}
\newcommand{\E}{\mathbb{E}}
\newcommand{\eps}{\varepsilon}
\newcommand{\dd}{\,\mathrm{d}}
\begin{document}
\section*{UnsolvedMath ID: 30002819}
\textbf{Source problem code:} catalogue entry
\subsection{Definitions and mathematical statement}
The supplied record contains two questions about nonconvex polytopes. Its first question does not specify the ambient dimension. The following formulation of that question uses the three-dimensional Euclidean interpretation suggested by the second question; this interpretation requires confirmation.

For a three-dimensional polyhedral solid \(P\subset\R^3\) with planar polygonal faces, let \(\alpha_P(F,v)\) denote the interior planar angle of a face \(F\) at one of its vertices \(v\), including a reflex angle when appropriate. Let \(\delta_P(e)\) denote the interior dihedral angle at an edge \(e\), measured on the side occupied by the solid, so that convex and reflex dihedral angles are distinguished. These definitions presuppose that the local face and solid geometry makes the indicated angles well defined. The source does not specify admissible degeneracies or its convention for dihedral angles.

Fix a combinatorial type \(K\), understood as the incidence structure of vertices, edges, and faces. For two realizations \(P,Q\) of \(K\), let \(\phi\) be an incidence-preserving correspondence. The first question asks whether, for every such pair of nonconvex realizations in the intended class,
\[
 \bigl(\delta_P(e)=\delta_Q(\phi(e))\text{ for every edge }e\bigr)
 \quad\Longrightarrow\quad
 \bigl(\alpha_P(F,v)=\alpha_Q(\phi(F),\phi(v))
 \text{ for every incident pair }(F,v)\bigr).
\]
The class of realizations, the ambient dimension, and any topological or nondegeneracy restrictions for this question remain unspecified in the supplied record.

For the second question, let \(P\subset\R^3\) be a nonconvex polyhedral solid whose boundary is homeomorphic to \(S^2\). For every face \(F\), write its containing plane as
\[
 H_F=\{x\in\R^3:\langle n_F,x\rangle=h_F\},
\]
where \(n_F\) is the unit normal pointing locally outward from the solid along the relative interior of \(F\). Its outward translated plane at distance \(\varepsilon\) is
\[
 H_F(\varepsilon)=H_F+\varepsilon n_F
 =\{x\in\R^3:\langle n_F,x\rangle=h_F+\varepsilon\}.
\]
Under the interpretation that an offset translates each face's containing plane while allowing its polygonal boundary to change, the second question is whether
\[
 \forall P\quad\exists\varepsilon_0(P)>0\quad
 \forall\varepsilon\in(0,\varepsilon_0(P))\quad
 \exists P_\varepsilon
\]
with \(P_\varepsilon\) an embedded polyhedral solid admitting a bijection \(F\mapsto F_\varepsilon\) from the faces of \(P\) to those of \(P_\varepsilon\), such that every \(F_\varepsilon\) has containing plane \(H_F(\varepsilon)\) and corresponding outward normal \(n_F\). The bijection expresses the requirement that every original face persists and no additional faces occur.

This offset formulation is conditional on the containing-plane interpretation. The supplied phrase ``each face translated outward'' could instead demand the literal equality \(F_\varepsilon=F+\varepsilon n_F\), which is a stronger, different requirement. The source also does not say whether the offset must preserve all incidences, retain spherical boundary, or converge to \(P\) in a specified sense. These unresolved requirements must be checked before either interpretation is adopted as the definitive statement.
\subsection{Short English statement}
The first question asks whether corresponding dihedral angles determine every angle within the faces of nonconvex polytopes with fixed combinatorics. Its dimension and admissible class require clarification. The second asks whether every nonconvex three-dimensional polyhedral solid with spherical boundary has an outward offset for every sufficiently small positive distance, retaining every face and introducing none. The meaning of translating a face, and the required topology and incidences of the offset, remain unresolved. The supplied record labels the combined problem partially solved; this is supplied dated curation, not an independently verified mathematical conclusion.
\subsection{Sources}
\begin{itemize}
\item Supplied local UnsolvedMath registry record, permanent ID \(30002819\), title: \emph{Dihedral Rigidity and Offsets of Nonconvex Polytopes}; exact source problem code: \texttt{catalogue entry}. Collection: Oberwolfach Reports Open Problems. The supplied proposed year is 2015; no individual proposer is identified.
\item Supplied source citation: \emph{Discrete Differential Geometry} (2015), Oberwolfach Reports, DOI 10.4171/owr/2015/13, pp.\ 661--729. \url{https://doi.org/10.4171/owr/2015/13}. The supplied background describes these reports as documenting mathematical workshops and their open-problem discussions. No precise problem page or individual attribution is supplied.
\item Supplied literature source: \emph{Straight skeletons and mitered offsets}. \url{https://doi.org/10.1007/s00454-016-9811-5}.
\item Supplied literature source: \emph{Volume conjecture and Stoker rigidity}. \url{https://arxiv.org/abs/2209.13439}.
\item Supplied literature source: \emph{OWR Report 13/2015}. \url{https://doi.org/10.4171/owr/2015/13}.
\item Literature assessment dated exactly \texttt{2026-08-21}: the supplied curation labels the problem \emph{partially solved}, states that offset-polytope existence has a general treatment and infinitesimal Stoker-type rigidity is known, and states that the full determination of face angles from dihedral angles for arbitrary nonconvex polytopes remains open. This assessment is supplied dated curation, not independently verified; its time-sensitive status, cited claims, hypotheses, and applicability to the two formulations require independent checking.
\item The supplied statement status is \texttt{corrected\_verified}. Its accompanying verification reports manual semantic, typographic, and LaTeX review of a meaning-preserving cleanup. This is a supplied editorial assertion, not independent verification of the mathematical formulation or literature status.
\end{itemize}
\end{document}
